% !TEX program = lualatex
\documentclass[12pt, a4paper]{article}

\usepackage{fontspec}
\setmainfont{TeX Gyre Pagella}

\usepackage{amsmath, amssymb, amsthm}
\usepackage{geometry}
\usepackage{microtype}
\usepackage{setspace}
\usepackage{titlesec}
\usepackage{hyperref}
\usepackage{enumitem}
\usepackage{booktabs}
\usepackage{graphicx}
\usepackage{epigraph}
\usepackage{mdframed}

\geometry{
  top=1.25in,
  bottom=1.25in,
  left=1.35in,
  right=1.35in
}

\onehalfspacing

\hypersetup{
  colorlinks=true,
  linkcolor=black,
  citecolor=black,
  urlcolor=black
}

% Theorem environments
\newtheoremstyle{named}
  {12pt}{12pt}{\itshape}{}{\bfseries}{.}{.5em}{}
\theoremstyle{named}
\newtheorem{theorem}{Theorem}
\newtheorem{definition}{Definition}
\newtheorem{corollary}{Corollary}
\newtheorem{proposition}{Proposition}

% Boxed theorem environment — theorem counter increments inside the box
\newmdenv[
  linewidth=0.8pt,
  innertopmargin=10pt,
  innerbottommargin=10pt,
  innerleftmargin=14pt,
  innerrightmargin=14pt,
  skipabove=14pt,
  skipbelow=14pt
]{thmbox}

\newenvironment{boxedtheorem}[1]{%
  \begin{thmbox}\begin{theorem}[#1]%
}{%
  \end{theorem}\end{thmbox}%
}

% Plain thmbox for non-theorem boxes (remarks, compressed statements)

\titleformat{\section}{\large\bfseries}{\thesection.}{0.5em}{}
\titleformat{\subsection}{\normalsize\bfseries}{\thesubsection.}{0.5em}{}

\title{%
  \textbf{When Systems Change Their Minds}\\[0.4em]
  \large Distributed Agency, Institutional Attractors,\\
  and the Reconfiguration of Economic Priorities
}
\author{Flyxion}
\date{June 2026}

% ---------------------------------------------------------------
\begin{document}

\maketitle
\thispagestyle{empty}

\begin{abstract}
Economic history is often narrated through great individuals, political
parties, or ideological movements. This essay proposes an alternative
interpretation. Rather than viewing economic transformations as the
direct product of centralized planners, we examine how large-scale
institutions function as distributed attention systems whose priority
structures evolve through selection, synchronization, and memory decay.
Through the transition from the Bretton Woods era to the contemporary
financialized economy, we investigate how changing priority structures
produce emergent forms of agency without requiring any single actor to
possess complete systemic control.

The essay develops a formal apparatus comprising one central definition
and five interconnected theorems describing the dynamics of institutional
transformation. The primary contribution
is the concept of \emph{Operational Fiber Death}: the condition in which
a system retains declarative knowledge of a prior configuration while
losing the procedural pathways required to reconstruct it. From this
definition, the paper derives theorems on reconstructability loss,
anticipatory collapse, approximate repair, opacity selection, and
self-concealment, culminating in two synthesis theorems on the dynamical
and ontological character of institutional change. The economic history
serves as empirical evidence for a broader claim: that large-scale
systems change principally by reconfiguring which futures remain
reachable, and that the deepest transformations are those which
eliminate not only the futures themselves but the capacity to recover
them.
\end{abstract}

\newpage
\tableofcontents
\newpage

% ---------------------------------------------------------------
\section{Introduction: The Problem of Large-Scale Change}

Historical explanation oscillates between two unsatisfactory poles. The
first attributes large-scale social transformations to impersonal
structural forces — the logic of capital, the pressures of geopolitics,
the imperatives of technology. The second attributes them to the
intentions of powerful individuals or coordinated elites: the memo
written, the rate set, the election won. Both approaches capture part of
the truth while obscuring a mechanism that neither can describe
adequately.

Systems do not merely change what they do. Over time, they change what
they remember how to do.

The loss of reconstructive capacity is difficult to observe precisely
because the processes that generate it simultaneously eliminate the
vantage points from which it could be recognized. Reachability closure
removes visible alternatives. Memory decay removes experienced
practitioners. Opacity selection removes auditability. By the time the
process completes, the observers best positioned to identify it have
themselves been selected out of the institutional environment.

The purpose of the present work is therefore not primarily historical
but diagnostic: to develop a formal language capable of describing how
societies lose access to previously reachable futures, and why that loss
tends to be recognized only retrospectively, after the procedural
pathways have already vanished.

\medskip

The choice of formalism requires justification. A skeptical reader might
fairly ask whether five theorems and an opacity functional add anything
that a careful narrative historian could not have said. The answer is
yes, but only once we recognize why narrative alone is insufficient:
narrative operates inside the same temporal embedding that the mechanism
exploits. A historian situated after the transition will naturally
describe the outcome as the product of intention, because the
alternatives have disappeared and the actors who could have testified
otherwise are gone. The formalism attempts — provisionally, fallibly —
to construct a vantage point outside that temporal structure. It does
not succeed perfectly. No formalism does. But the attempt itself has
diagnostic value, because it forces explicit statement of what is being
claimed about mechanism, timing, and irreversibility — claims that
narrative can make only implicitly, and therefore cannot easily be
examined or contested.

\medskip

The paper proceeds as follows. Section~2 introduces the background
concepts of agency, attractors, and distributed optimization.
Section~3 describes the Bretton Woods configuration and its internal
logic. Sections~4 and~5 trace the crisis, synchronization, and
transition dynamics. Section~6 introduces the formal apparatus,
opening with the definition of Operational Fiber Death and proceeding
through Theorems~1--4 in logical dependency order. Section~7 develops
the opacity and financialization analysis. Section~8 states Theorem~5
on self-concealment and the Emergent Agency synthesis. Section~9
addresses repair. Section~10 generalizes the framework beyond the
economic case. The conclusion states the two synthesis theorems and
returns to the essay's central proposition.

% ---------------------------------------------------------------
\section{Agency Without an Agent}

The human tendency to anthropomorphize complex systems is not merely a
cognitive error. It is a reasonable inference under conditions of
opacity. When a system produces coherent, persistent, directional
behavior, the most parsimonious explanation available to an observer
without access to the underlying mechanism is that the system has
purposes. The inference is wrong about the mechanism but correct about
the behavior.

A thermostat does not want the room to be warm. An ant colony does not
decide to build chambers. A market does not choose to clear. A
bureaucracy does not intend to expand. Yet all of these systems produce
behavior that, observed from the outside, resembles intention. The
source of that resemblance is not hidden consciousness but recursive
self-correction: feedback loops that cause the system to return toward
certain states and away from others, producing apparent goals without
requiring any goal-bearing agent.

The concept of an \emph{attractor} formalizes this observation. In
dynamical systems theory, an attractor is a region of state space toward
which trajectories converge under the system's evolution equations
\cite{strogatz2015,wiener1948,meadows2008}. A system with a stable attractor will
return to that region after perturbation. It will appear to \emph{seek}
the attractor. The appearance is real in the sense that the behavior is
reliable; the seeking is illusory in the sense that no seeker exists.

Large institutional systems — economies, governments, legal orders,
scientific communities — possess attractors in this sense
\cite{simon1962,holland1992,holland1995}. The behavior of such a system over time
is not random, and it is not fully controlled by any individual actor.
It is shaped by the distribution of incentives, feedback mechanisms,
selection pressures, and constraint fields that constitute the
institution's effective priority structure \cite{douglassnorth1990,
marcholsen1989}. Understanding how that structure changes is the central
problem this essay addresses.

% ---------------------------------------------------------------
\section{The Bretton Woods Configuration}

The postwar economic order established at Bretton Woods in 1944
represented a specific institutional configuration with a specific
priority vector \cite{polanyi1944,galbraith1967}. Fixed exchange rates
constrained competitive devaluation. Capital controls limited
cross-border financial flows. Strong labor unions maintained wage
bargaining power. Progressive taxation compressed income distributions.
Public investment in infrastructure, housing, and education sustained
broad-based growth.

The system's dominant objective function — to use the language we will
formalize below — weighted employment, social cohesion, reconstruction,
and political stability heavily \cite{minsky1986}. Financial returns
existed within this framework but were subordinate to it. Capital was a
means; production and employment were ends.

This configuration was not the product of a unified plan. It emerged
from the particular balance of political forces, international pressures,
and institutional memories of the 1930s that shaped postwar settlements
across the industrial democracies \cite{polanyi1944,streeck2016}. Its
coherence was real but distributed: different institutions in different
countries pursued different local objectives, yet the aggregate behavior
of the system was recognizably consistent \cite{rodrik2011}.

The Bretton Woods configuration sustained three decades of relatively
broad-based growth, falling inequality, and expanding public
infrastructure across the industrial world \cite{piketty2014}. It also
accumulated the contradictions that would eventually destabilize it:
dollar outflows from military spending and corporate investment abroad,
inflationary pressures from full-employment commitments, rising commodity
prices, and the increasing mobility of capital as communications
technology reduced transaction costs \cite{harvey2005,dumenil2004}.

% ---------------------------------------------------------------
\section{Crisis as Attractor Destabilization}

\subsection{The Formal Framework}

The following formalism should be understood as a dynamical model of
institutional evolution rather than a derivation from economic first
principles. It supplies a precise language for mechanisms the historical
narrative identifies but cannot itself formalize.

Let $x(t)$ denote the state of the economy at time $t$, encompassing
wages, unemployment, inflation, asset prices, capital mobility, public
spending, and labor power. Let $\mathbf{p}_i(t) \in \mathbb{R}^n$
denote the priority vector of institution $I_i$, where each coordinate
represents a weighted institutional objective. The aggregate
institutional field is:
\[
  P(t) = \sum_{i=1}^{N} w_i \mathbf{p}_i(t),
\]
where $w_i$ denotes the relative power and amplification capacity of
institution $I_i$.

Define an institutional potential function:
\[
  V(x,t) = -\langle P(t),\, F(x)\rangle,
\]
where $F(x)$ maps the economic state into institutionally visible
features. The system evolves approximately by gradient descent:
\[
  \frac{dx}{dt} = -\nabla_x V(x,t) + \xi(t),
\]
where $\xi(t)$ represents exogenous shocks: oil price movements,
political conflict, technological change, and geopolitical disruption.

\subsection{Instability}

A crisis occurs when the existing institutional configuration can no
longer minimize its own potential function. The old attractor
$x^{\ast}$ becomes unstable when:
\[
  \nabla_x V(x^{\ast},t) = 0
\]
but the Hessian loses positive definiteness:
\[
  \nabla_x^2 V(x^{\ast},t) \not\succ 0.
\]

Constraint accumulation in the years before a visible transition can be
interpreted as the institutional equivalent of entropy-density growth:
contradictory objectives, rising cost pressures, and failing legitimacy
accumulate before any observable phase change occurs. The Hessian
instability criterion supplies the local stability condition that marks
when accumulated tension exceeds the restoring capacity of the existing
configuration \cite{strogatz2015,turchin2016}.

\subsection{The Empirical Transition}

The destabilizing forces of the early 1970s are well documented
\cite{harvey2005,minsky1986,krippner2011}: the end of dollar-gold
convertibility in 1971, the oil shocks of 1973 and 1979, the resulting
stagflation that conventional Keynesian tools could not resolve, and
the declining confidence in the postwar economic order that followed.
These events do not collectively constitute a plan. They constitute a
fitness landscape deformation: the institutional configurations that had
reproduced successfully under Bretton Woods began to fail, and the space
of viable alternatives temporarily broadened \cite{turchin2016}.

% ---------------------------------------------------------------
\section{Synchronization and Transition}

\subsection{The Replicator Dynamics}

The standard account of the neoliberal transition emphasizes ideology:
the Powell Memo of 1971 \cite{powell1971}, the Mont Pelerin Society,
the Chicago School, the think tanks that proliferated through the 1970s
\cite{mirowski2013}. This account is not wrong, but it is incomplete.
It treats priority vector rotation as the primary mechanism:
\[
  \mathbf{p}_i(t) \to \mathbf{p}_i(t + \Delta t).
\]
But institutions often do not change people. They change who survives.

A more precise formulation replaces individual vector rotation with
population dynamics. Let $\rho(\mathbf{p}, t)$ denote the density of
institutional actors possessing priority vector $\mathbf{p}$ at time
$t$, and let $F(\mathbf{p}, t)$ denote the institutional fitness of
that vector. The evolution of the population follows:
\[
  \frac{\partial \rho(\mathbf{p},t)}{\partial t}
  = \rho(\mathbf{p},t)\bigl(F(\mathbf{p},t) - \bar{F}(t)\bigr),
\]
where:
\[
  \bar{F}(t) = \int F(\mathbf{p},t)\,\rho(\mathbf{p},t)\,d\mathbf{p}.
\]

This is a replicator equation \cite{nelsonwinter1982,march1991}. The
equation is borrowed from evolutionary biology and applied here in the
tradition established by Nelson and Winter \cite{nelsonwinter1982}, who
explicitly adopt formal analogy as a methodological tool: the replicator
dynamics are not claimed to be literally true of promotion committees but
to capture, at the population level, the net effect of differential
retention across a large number of local hiring and advancement decisions.
Actors whose priority vectors match the emerging institutional fitness
landscape are promoted, retained, and hired. Actors whose vectors do not
match are marginalized, passed over, or retired. The institution therefore
evolves without anyone necessarily changing their beliefs. The population
composition changes; the belief distribution of any individual actor may
remain constant throughout \cite{cyertmarch1963,weick1995}.

The fitness functional itself decomposes into multiple components:
\[
  F(\mathbf{p},t)
  = F_{\text{economic}} + F_{\text{institutional}}
  + F_{\text{cultural}} + F_{\text{network}}.
\]
During the crisis period, $F$ deforms: vectors that were previously
adaptive become maladaptive, and vectors that were previously marginal
become viable. This deformation is the window during which
synchronization mechanisms matter — not because they persuade
individuals, but because they arrive at the moment when the selection
environment is briefly underdetermined.

\subsection{Synchronization}

Let institutional alignment be measured by:
\[
  A(t) = \frac{1}{N(N-1)}
  \sum_{i \neq j}
  \frac{\langle \mathbf{p}_i(t), \mathbf{p}_j(t)\rangle}
       {|\mathbf{p}_i(t)||\mathbf{p}_j(t)|}.
\]
When $A(t)$ is low, institutions pull in competing directions. When
$A(t)$ rises, the system begins to behave as if it possessed a unified
will.

The Powell Memo \cite{powell1971}, think tanks, industry associations,
and elite policy networks function as synchronization mechanisms
\cite{mirowski2013,harvey2005}. They do not create coherence from
nothing. They align already-existing institutional vectors into a more
coherent field at the moment when fitness landscape deformation has made
alignment advantageous. Timing, not content, is what made such
interventions effective.

\subsection{The Volcker Shock as Control Input}

Let $\pi(t)$ denote inflation, $u(t)$ unemployment, and $L(t)$ labor
bargaining power. A simplified monetary control rule is:
\[
  r(t) = r_0 + \alpha(\pi(t) - \pi^{\ast}),
\]
where $r(t)$ is the interest rate and $\pi^{\ast}$ is the target
inflation. The labor channel evolves as:
\[
  \frac{dL}{dt} = -\beta u(t) - \gamma C(t),
\]
where $C(t)$ represents deindustrialization, plant closures, and strike
fragility. Since high interest rates suppress investment and raise
unemployment:
\[
  r(t)\uparrow \;\Rightarrow\; u(t)\uparrow \;\Rightarrow\; L(t)\downarrow.
\]
The Federal Reserve's rate increases to near 20\% in the early 1980s
were therefore not merely an anti-inflation measure. They were a control
input that restructured the fitness landscape for labor-aligned priority
vectors.

The language of control inputs requires a clarification that applies
throughout this paper. Describing Volcker's rate decisions, Thatcher's
union legislation, or the Powell Memo as inputs to an attractor system
does not deny that these actors had intentions. It distinguishes two
levels of agency that the standard historical narrative conflates.
\emph{Local intentionality} refers to specific actors pursuing specific
goals within a fitness landscape they did not design and largely did not
understand as a system. \emph{Global apparent intentionality} refers to
the coherent long-run trajectory of the system as a whole — the sense
in which the economy appears to have \emph{decided} to financialize, or
the state appears to have \emph{chosen} to withdraw from labor
protection. The paper's claim is about the second level, not the first.
Volcker intended to break inflation; he did not intend, and could not
have designed, the full reachability reconfiguration that followed from
his decision interacting with the existing fitness landscape, the
synchronization already underway, and the global debt mechanics that
propagated the shock. Local intentions are real. Global apparent
intention is emergent. The error — in conspiracy thinking and in
ideological history alike — is to read the second as evidence of a
scaled-up version of the first.

\subsection{Global Propagation}

Let a developing country hold dollar-denominated debt $D$. Its debt
service burden is $B(t) = r_{\mathrm{USD}}(t)\,D$, so:
\[
  \frac{dB}{dr_{\mathrm{USD}}} = D > 0.
\]
A rise in U.S. interest rates mechanically increases repayment pressure.
When $B(t) > Y(t) - G_{\min}(t)$, where $Y(t)$ is national income and
$G_{\min}(t)$ is the minimum politically sustainable public expenditure,
the country enters debt crisis \cite{stiglitz2002,rodrik2011}. The
resulting IMF structural adjustment programs — privatization, austerity,
deregulation, labor market liberalization — propagated the priority
vector rotation globally through a financial mechanism that required no
ideological conversion \cite{harvey2005,dumenil2004}. The fitness
landscape of recipient economies was simply restructured by the terms
attached to necessary credit.

% ---------------------------------------------------------------
\section{The Formal Apparatus}

We now state the paper's primary formal contributions in their logical
dependency order.

\subsection{Definition: Operational Fiber Death}

\begin{definition}[Operational Fiber Death]
Let $\pi: X \to M$ be a projection from institutional state space $X$
onto the publicly observable manifold $M$. A state $m \in M$ exhibits
\emph{operational fiber death} when:
\[
  \pi^{-1}(m) \neq \varnothing
\]
but the measure of admissible reconstruction pathways $\mathcal{P}_m$
satisfies:
\[
  \mu(\mathcal{P}_m) \to 0.
\]
That is: the preimage of $m$ remains formally non-empty, but the
procedural operators required to construct or reach elements of that
preimage have decayed to measure zero.
\end{definition}

Operational Fiber Death is strictly stronger than information loss,
stronger than forgetting, and stronger than political infeasibility.

\begin{thmbox}
The possibility still exists in principle. The route has disappeared
in practice.
\end{thmbox}

To make this concrete in the institutional context: a society may
retain complete declarative knowledge that capital controls existed,
that industrial unions coordinated wage bargaining, that public housing
regimes allocated residential space, and that investment banks were
legally separated from commercial banks. It may possess historical
archives, academic literature, and policy documents describing how each
of these operated. Knowledge remains. The procedural capacity — the
trained personnel, established protocols, regulatory expertise, and
operational infrastructure — has decayed to the point where
reconstruction would require building an entirely new institutional
apparatus rather than restoring an existing one.

\subsection{Theorem 1: Reconstructability Loss}

\begin{boxedtheorem}{Reconstructability Loss}
Let $\mathcal{A}(t)$ denote the admissible set of reachable social
futures at time $t$. Reachability closure $\mathrm{Vol}(\mathcal{A}_{\text{labor}}) \to 0$
is not a temporary equilibrium displacement. Once the replicator
dynamics have operated for a sufficient period, the closure becomes
structural: the institutional population required to operate the
closed configuration has been selected out of the system. Reachability
closure is reversible in principle; reconstructability loss is
irreversible without procedural memory. Formally,
$\mathcal{A}_{\text{old}} \not\subset \mathcal{A}(t)$ establishes that
the old trajectory is no longer reachable under current constraints,
while $A_p(t) \to 0$ establishes that the system no longer possesses
the procedural capacity to reconstruct the operators that once made
that trajectory available. These are distinct conditions. The second
is strictly stronger.
\end{boxedtheorem}

The mechanism operates through two channels. First, actors with
non-aligned priority vectors are not promoted and retire without
successors trained in their orientation \cite{march1991,teece1997}.
Second, institutional memory decays at different rates for different
knowledge types \cite{nelsonwinter1982,weick1995}. Let $A(t)$
decompose into declarative memory $A_d$ and procedural memory $A_p$:
\[
  \frac{dA_d}{dt} = -\lambda_d D(t), \qquad
  \frac{dA_p}{dt} = -\lambda_p D(t),
\]
where $D(t)$ represents demographic turnover and $\lambda_p \gg
\lambda_d$. Procedural memory decays faster because it requires
continuous practice to maintain \cite{arrow1962,arthur1994}. After
approximately one generational cycle, $A_p \to 0$ while $A_d$ remains
substantial: people know a configuration existed without retaining the
capacity to instantiate it. This is operational fiber death at the
institutional scale.



\subsection{Theorem 2: Anticipatory Collapse}

\begin{theorem}[Anticipatory Collapse]
The contraction of the admissible set begins before the institutional
transition completes. Formally, there exists $t^{\ast} < t_{\text{transition}}$
such that:
\[
  \frac{d}{dt}\,\mathrm{Vol}(\mathcal{A}_{\text{labor}})\big|_{t=t^{\ast}} < 0,
\]
that is, the admissible set is already shrinking during the crisis and
synchronization phases, before the new attractor has stabilized.
\end{theorem}

This theorem has an important corollary for the interpretation of
political resistance. At any moment during the transition, the actors
most aware of what is being lost are also the actors whose procedural
knowledge is actively decaying. Their knowledge of the prior
configuration is real, but their capacity to articulate what
specifically is being foreclosed diminishes precisely as the foreclosure
accelerates. This temporal asymmetry explains why resistance to
large-scale institutional transitions tends to be described, by
historians writing afterward, as reactive and nostalgic rather than
prescient — even when the resistors correctly identified the mechanism
in real time.

\subsection{Theorem 3: Approximate Repair}

\begin{theorem}[Approximate Repair]
Under operational fiber death, exact institutional restoration is
unavailable. The maximum achievable repair is functional recovery:
\[
  R_{\epsilon}(S_{\text{broken}})
  \approx_{\mathrm{function}} S_{\text{old}},
  \qquad
  R_{\epsilon}(S_{\text{broken}})
  \not\cong S_{\text{old}}.
\]
A repaired institution may reproduce the lost function — broad
bargaining power, capital discipline, housing access, public investment
— without reproducing the historical machinery that once achieved it.
\end{theorem}

This theorem is a corollary of Theorem~1 rather than an independent
result. It is stated separately because it has direct implications for
political discourse. Most reform proposals implicitly assume that repair
is restoration: that the correct objective is to rebuild the institutions
that existed prior to the transition \cite{streeck2016,mirowski2013}.
The theorem implies that this framing is not merely politically difficult
but formally unavailable once procedural memory has decayed. What
remains possible is reconstruction that achieves functional equivalence
through new institutional forms \cite{ostrom1990,williamson1985}. This
is not a counsel of despair; it is a correction of the target.

The reachability framework formalizes the distinction between
distributional and topological accounts of inequality
\cite{piketty2014,granovetter1985}. Most standard treatments operate
on a distributional ontology:
\[
  \text{Inequality} = \text{Unequal allocation.}
\]
The present framework implies a different ontology:
\[
  \text{Inequality} = \text{Unequal reachability.}
\]
Suppose each citizen $i$ possesses a future opportunity set
$\mathcal{A}_i$. Define aggregate social freedom as:
\[
  \mathcal{F} = \sum_i \mathrm{Vol}(\mathcal{A}_i).
\]
A society can increase aggregate GDP while reducing $\mathcal{F}$.
A housing market can appreciate in value while reducing
$\mathrm{Vol}(\mathcal{A}_{\text{young}})$. An educational system can
expand expenditure while reducing
$\mathrm{Vol}(\mathcal{A}_{\text{student}})$. A labor market can
increase productivity while reducing
$\mathrm{Vol}(\mathcal{A}_{\text{worker}})$. This is a substantially
deeper claim than conventional inequality metrics capture.

\begin{proposition}[Reachability Principle]
The effective freedom of an actor is proportional not to the resources
possessed but to the volume of admissible futures accessible from their
current state \cite{sen1999}:
\[
  \text{Freedom}_i \;\propto\; \mathrm{Vol}(\mathcal{A}_i).
\]
Aggregate social freedom is therefore:
\[
  \mathcal{F} = \sum_i \mathrm{Vol}(\mathcal{A}_i),
\]
and inequality is best measured not by distributional distance from the
mean but by the variance of $\mathrm{Vol}(\mathcal{A}_i)$ across the
population.
\end{proposition}

% ---------------------------------------------------------------
\section{Opacity, Financialization, and Self-Concealment}

\subsection{The Opacity Operator}

Define a visibility operator:
\[
  \pi: X \to M,
\]
mapping the true institutional state space $X$ onto the publicly
observable manifold $M$. Opacity is information loss under this
projection:
\[
  \Omega = H(X) - H(\pi(X)).
\]
When $\Omega$ is small, participants can trace causal chains and
understand local mechanisms. When $\Omega$ becomes large, the system
appears to possess unified motives because the underlying interactions
are inaccessible.

Opacity decomposes into two qualitatively distinct components:
\[
  \Omega = \Omega_S + \Omega_M,
\]
where $\Omega_S$ is \emph{structural opacity} arising from genuine
system depth, scale, and recursion, and $\Omega_M$ is
\emph{manufactured opacity} arising from strategic concealment,
regulatory arbitrage, and deliberate complexity. Only $\Omega_M$ enters
directly into institutional optimization:
\[
  F(\mathbf{p},t) = F_0(\mathbf{p},t) + \kappa\,\Omega_M.
\]
This implies that some systems actively select for actors capable of
increasing opacity — not through explicit coordination, but because
increased opacity raises the fitness of the actors who generate it.

\subsection{Theorem 4: Opacity Selection}

\begin{theorem}[Opacity Selection]
Institutional measurement degrades as the mechanism operates.
Specifically, the component $\Omega_M$ of institutional opacity is
itself a fitness-enhancing trait under conditions of reachability
closure and procedural memory decay. Systems will therefore evolve
toward increasing manufactured opacity through selection, independently
of any deliberate strategy to obstruct accountability.
\end{theorem}

The empirical prediction is not that opacity will increase without
limit. The fitness functional is concave in $\Omega_M$:
\[
  F(\Omega_M) = F_0 + \kappa\,\Omega_M - \eta\,\Omega_M^2.
\]
The optimum is:
\[
  \Omega_M^{\ast} = \frac{\kappa}{2\eta}.
\]

\begin{proposition}[Optimal Opacity]
Institutions selected for opacity will converge toward a locally
optimal opacity level $\Omega_M^{\ast}$ rather than maximal opacity.
\end{proposition}

\begin{corollary}[Systemic Fragility]
When $\Omega_M > \Omega_M^{\ast}$, internal coordination costs
dominate external accountability benefits. The institution becomes
fragile: it is protected from external scrutiny but unable to
understand its own risk structure.
\end{corollary}

The financial crisis of 2008 illustrates this corollary
\cite{minsky1986,taleb2012,krippner2011}. The same opacity mechanisms
— securitization chains, off-balance-sheet vehicles, derivatives
complexity, algorithmic trading — that protected financial actors from
regulatory oversight eventually prevented the institutions themselves
from accurately modeling their own exposures. Opacity ceased to be
protective and became destabilizing. This is exactly what a concave
fitness function with interior optimum predicts.

A note on empirical scope is warranted here. The parameters $\kappa$,
$\eta$, and $\Omega_M$ are not directly measurable in any existing
institutional dataset, and the paper does not claim otherwise. The
fitness functional is a modeling schema rather than an empirically
calibrated law. What it offers is a qualitative, directional prediction:
institutions under opacity selection will exhibit internal coordination
failure before reaching maximal opacity, and the signature of that
failure — risk blindness coexisting with external impermeability — is
observable without measuring the parameters themselves. Calibration is a
research program that the present framework makes possible to specify;
it is not a precondition for the framework's validity. The emergence
threshold $\Omega > \Omega_{\text{threshold}}$ in the Emergent Agency
theorem operates differently: it is a conceptual marker identifying the
condition under which alignment produces apparent intentionality, not an
independently measurable criterion. These two uses of opacity should be
kept distinct.

Financialization can now be described formally as a signal reweighting
\cite{krippner2011,dumenil2004}. Let economic output decompose into
productive income $Y_P$ and financial income $Y_F$:
\[
  Y(t) = Y_P(t) + Y_F(t).
\]
Financialization occurs when institutional reward functions increasingly
privilege $Y_F$ over $Y_P$:
\[
  U_{\text{firm}} = a\,Y_P + b\,Y_F, \quad b \gg a.
\]
Under this regime, firms optimize asset appreciation, shareholder
returns, buybacks, leverage, and speculative instruments rather than
wages, durability, or productive capacity. The economy becomes, in
effect, a machine for processing financial information — one in which
the signals that drive institutional behavior have become increasingly
abstracted from the material processes those institutions nominally
govern.

% ---------------------------------------------------------------
\section{Apparent Agency}

\subsection{Theorem 5: Self-Concealment}

\begin{boxedtheorem}{Self-Concealment}
Let $V(t)$ denote the volume of accessible diagnostic vantage points
from which the transition mechanism can be identified. Throughout the
process of reachability closure, procedural memory decay, and opacity
selection:
\[
  \frac{dV}{dt} < 0.
\]
The mechanism progressively eliminates the vantage points from which
it could be diagnosed. $V(t) \to 0$ is not a contingent outcome but a
structural feature of the dependency chain:
\[
\begin{aligned}
  &\text{Closure}
  \;\to\; \text{Memory Decay} \\
  &\quad\to\; \text{Opacity Selection}
  \;\to\; \text{Apparent Agency.}
\end{aligned}
\]
\end{boxedtheorem}

The logic is sequential. Reachability closure removes the visible
alternatives that would make the closure recognizable as a loss.
Procedural memory decay removes the experienced practitioners who could
articulate what was being foreclosed. Opacity selection removes the
audit capacity that would make the selection mechanism legible. At each
stage, the set of available diagnostic positions shrinks. By the time
the process is complete, the observers best positioned to identify it
have been removed from the institutional environment.

This theorem explains why institutional transformations are almost
always recognized retrospectively. The recognition feels like
historical understanding; it is in fact confirmation of operational
fiber death. It also explains why the present formalism is necessary:
narrative history operates inside the temporal embedding that the
mechanism exploits, and formalization — however provisional — creates
explicit statements that can be examined independently of the narrative
momentum that carries historical description toward intentional
attributions.

\subsection{Emergent Agency (Synthesis)}

\begin{theorem}[Emergent Agency]
A distributed institutional system exhibits apparent agency when the
following conditions hold simultaneously:
\[
  \text{Persistence} + \text{Feedback} + \text{Alignment}
  + \text{Opacity}
  \;\Longrightarrow\; \text{Apparent Agency.}
\]
More precisely: if $A(t) > \theta$ for some alignment threshold
$\theta$, if the system recursively updates according to:
\[
  \mathbf{p}_i(t+1)
  = \mathbf{p}_i(t)
  + \eta\,\nabla_{\mathbf{p}_i} R_i(x(t)),
\]
and if $\Omega > \Omega_{\text{threshold}}$, then the aggregate field
$P(t)$ generates coherent macroscopic trajectories even when no
individual actor controls the whole.
\end{theorem}

The role of opacity in this theorem distinguishes it from standard
accounts of emergent coordination. Without opacity, aligned institutions
are merely visible coordination: observers can see the mechanism and
attribute behavior to distributed optimization. With opacity, the same
alignment appears as unified intention. The phenomenology of conspiracy
thinking is precisely the observer's correct detection of alignment
through an opacity screen that prevents them from seeing the distributed
mechanism beneath. The conspiracy theorist is often right about the
behavior and wrong about the architecture.

\newpage

The dependency chain can now be stated in its complete form:

\begin{thmbox}
\[
\begin{aligned}
  &\text{Reachability Closure}
  \;\to\; \text{Procedural Memory Decay} \\
  &\quad\to\; \text{Opacity Selection}
  \;\to\; \text{Apparent Agency.}
\end{aligned}
\]
Each stage removes explanatory alternatives. First, alternative
trajectories become inaccessible. Second, the operational knowledge
required to enact them disappears. Third, opacity becomes adaptive
because comparison classes have vanished. Fourth, observers attribute
unified intention to what remains.
\end{thmbox}

% ---------------------------------------------------------------
\section{Repair and Reconstruction}

The foregoing framework implies a specific account of institutional
repair that differs substantially from both conservative restoration
and progressive replacement.

Repair, as a formal category, is distinct from replacement
\cite{marcholsen1989,williamson1985}. Replacement destroys the existing
configuration and substitutes a new one, losing whatever institutional
memory the existing configuration embodies. Repair attempts to recover
function while preserving structural continuity. The distinction matters
because institutional memory — including memory of how to navigate
regulatory structures, negotiate with counterparties, and maintain
operational coherence under stress — is itself a productive resource
\cite{hayek1945,nelsonwinter1982} that is destroyed by replacement and
potentially recoverable through repair.

However, Theorem~1 establishes that exact restoration is unavailable
once operational fiber death has occurred: $A_p(t) \to 0$ means the
procedural operators for the old configuration have decayed beyond
recovery through ordinary institutional continuity. The repair that
remains possible is therefore approximate: functional recovery without
formal identity. Institutions designed to recover functions lost through
procedural decay cannot be designed by analogy to the historical
institutions that originally served those functions, because the
trained personnel, regulatory environment, and counterpart institutions
those historical forms required no longer exist. A functional equivalent
must be designed for the current environment.

This observation has a second practical implication. The policy
discourse that frames institutional reform as a choice between
nostalgic restoration and radical replacement is falsely dichotomous
\cite{ostrom1990,page2018}. Theorem~3 implies a third category:
functional reconstruction under memory loss. This is the appropriate
target for institutional design once the extent of procedural decay has
been honestly assessed.

% ---------------------------------------------------------------
\section{Generalization Beyond Political Economy}

The framework developed in the preceding sections applies to any
adaptive system that undergoes reachability closure, exhibits
differential memory decay rates, generates opacity through selection,
and thereby produces apparent agency. The economic case is illustrative,
not definitional.

Scientific paradigms undergo analogous transitions \cite{kuhn1962}.
Experimental techniques fall out of use. Instrument-building skills
decay. Theoretical frameworks that were once pursued as live options
become effectively unrecoverable because the training pipelines that
produced practitioners have been discontinued. A paradigm shift in
Kuhn's sense \cite{kuhn1962} is partly a reachability reconfiguration:
after sufficient time, the prior paradigm is not merely unfashionable
but operationally inaccessible — a case of operational fiber death at
the level of scientific practice rather than economic institution.

Software ecosystems exhibit the same dynamics. Programming languages,
architectural patterns, and system designs fall out of production use.
Organizations that once possessed the capacity to modify or extend
certain systems eventually retain only the declarative knowledge that
those systems exist; $A_p \to 0$ while $A_d$ persists in documentation
that no one can execute.

Religious traditions, legal systems, and educational institutions all
exhibit operational fiber death over sufficiently long time horizons.
The capacity to perform certain liturgical functions, argue certain
legal forms, or teach certain subjects requires continuous practice
maintained across generations. When the chain breaks — through
persecution, institutional collapse, or sustained disuse — the
knowledge of what was done may survive in texts while the knowledge of
how to do it disappears.

The generalization to artificial systems is particularly relevant. A
large-scale machine learning system optimizing an implicit objective
function through feedback on outcomes is a distributed institutional
system in exactly the sense formalized here. It exhibits persistence,
feedback, and alignment. If its internal operations are opaque to
external observers — and if the selection mechanisms that shaped its
optimization are not legible from its outputs — it will exhibit
apparent agency. The question this essay has been asking throughout —
\emph{what is the system actually optimizing, and who can see inside
it?} — is precisely the question the alignment research community is
attempting to answer for artificial systems.

% ---------------------------------------------------------------
\section{Toward a Theory of Institutional Attractors}

The economic history examined in this essay demonstrates the mechanism.
The formal apparatus developed in response to it claims a much wider
jurisdiction. The preceding theorem chain established the dependency
relations among its components:
\[
\begin{aligned}
  &\text{OFD}
  \;\to\; \text{Reconstructability Loss}
  \;\to\; \text{Anticipatory Collapse} \\
  &\quad\to\; \text{Opacity Selection}
  \;\to\; \text{Self-Concealment.}
\end{aligned}
\]
We are now in a position to state the two synthesis theorems of which
this chain is the derivation.

\begin{thmbox}
\textbf{Theorem (Dynamical).} \emph{Large systems change principally by
reconfiguring which futures remain reachable.}
\[
  P_{\text{Bretton}}
  \to P_{\text{Crisis}}
  \to P_{\text{Monetarist}}
  \to P_{\text{Financialized}}.
\]
Each transition is not merely a redistribution of resources within a
fixed possibility space. It is a transformation of the space itself:
\[
  \mathcal{A}(t+1) = \mathcal{T}_{\text{policy}}(\mathcal{A}(t)).
\]
\end{thmbox}

\begin{thmbox}
\textbf{Theorem (Ontological).} \emph{System-level intention can emerge
from recursive alignment long before anything resembling a mind
appears.}

The apparent agency of large institutions — the sense in which markets
appear to want things, governments appear to decide things, and
corporations appear to believe things — is not an illusion in the sense
of being false, nor is it a reality in the sense of implying a
subject. It is a structural feature of distributed systems that are
sufficiently aligned, sufficiently persistent, sufficiently recursive,
and sufficiently opaque.
\end{thmbox}

The first theorem is demonstrated by the historical case. The second is
implied by the first: once we see how priority vector alignment produces
coherent macroscopic trajectories without centralized control, we see
why those trajectories appear intentional to observers who lack access
to the distributed mechanism.

The task of political repair is therefore not merely to replace leaders
or redistribute resources. It is to alter the priority vectors, fitness
landscapes, feedback channels, and admissibility constraints through
which institutions determine what futures remain reachable. This is a
substantially harder problem than electoral politics addresses, and a
substantially more tractable problem than revolutionary replacement
proposes. It is the problem of institutional redesign under conditions
of partial memory loss — reconstruction that recovers function without
recovering form.

% ---------------------------------------------------------------
\section*{Conclusion}

The central question is not whether a hidden group controls society.
The more important question is: what objective functions are our
institutions currently optimizing, and what futures does that
optimization make unreachable?

The present essay has argued that this question cannot be answered
through historical narrative alone, because narrative operates inside
the same temporal embedding that the mechanism exploits. It has
proposed a formal apparatus — priority vectors, replicator dynamics,
operational fiber death, opacity selection, approximate repair — as an
attempt to construct, provisionally, a vantage point from which the
mechanism can be described without being subject to its
self-concealment.

The remaining gaps in that apparatus identify areas where the framework
has not yet been formally specified and therefore indicate directions for
future work: empirical calibration of the opacity functional, micro-level
derivation of the replicator dynamics from promotion and hiring
mechanisms, and extension of the reachability formalism to domains
beyond political economy.

The practical implication follows directly from Theorem~3. The
appropriate political target is not restoration of prior institutional
forms — Theorem~1 establishes that this is formally unavailable once
procedural memory has decayed — but functional reconstruction: the
design of new institutional forms capable of recovering the functions
that were lost. Stable employment, compressed inequality, capital
discipline, and broad-based investment are functions, not forms. They
were once achieved by one set of institutional arrangements; they could
in principle be achieved by others that have never yet existed. The
question is what those arrangements would need to look like in an
environment where the old procedural knowledge is gone and the old
counterpart institutions no longer exist. That question has not been
seriously asked, because the dominant reform discourse has been
organized around restoration rather than reconstruction. This paper
offers a reason to ask it instead.

The behavior of large-scale systems will continue to appear mysterious,
intentional, and conspiratorial as long as the underlying mechanism
remains invisible. The underlying mechanism remains invisible in part
because it progressively eliminates the vantage points from which it
could be seen. The appropriate response is not to conclude that no
mechanism exists, but to build the formal language required to describe
one that systematically resists description.

\bigskip

Institutions do not merely constrain what societies can do. Over time,
they alter what societies remember how to do. The deepest political
transformations therefore occur not when resources are redistributed,
but when reconstructive pathways disappear.

A civilization loses a future twice: first when it ceases to pursue it,
and later when it forgets how it was ever reached.


\newpage 
\section*{Appendices}

\appendix

\section{A. Reachability Geometry}

The central object of the paper is the admissible future set
$\mathcal{A}(t)$. We may define a reachability metric on institutional
state space by

\[
d_R(x,y)
=
\inf_{\gamma}
\int_0^1
c(\gamma(s))
\, ds,
\]

where $\gamma$ ranges over admissible trajectories connecting $x$ and
$y$, and $c$ denotes the institutional constraint density.

The reachable future volume from state $x$ is

\[
\mathcal{V}(x)
=
\int_{\mathcal{A}(x)}
d\mu.
\]

The local contraction rate is therefore

\[
\Lambda(x,t)
=
-
\frac{d}{dt}
\log \mathcal{V}(x,t).
\]

Large positive values of $\Lambda$ correspond to rapid closure of
institutional possibilities.

\begin{proposition}
If $\Lambda(x,t) > 0$ over an interval
$[t_0,t_1]$, then the admissible future volume satisfies

\[
\mathcal{V}(t_1)
=
\mathcal{V}(t_0)
\exp
\left(
-\int_{t_0}^{t_1}
\Lambda(\tau)\,d\tau
\right).
\]

Hence persistent closure produces exponential loss of reachable
futures.
\end{proposition}

\section{B. Procedural Memory Dynamics}

Let procedural memory be represented by a vector

\[
\mathbf{m}(t)
=
(m_1,\ldots,m_n).
\]

Each component corresponds to an institutional capability:
industrial policy, union negotiation, capital control management,
public housing administration, and so forth.

Assume

\[
\frac{dm_i}{dt}
=
-\lambda_i m_i
+
u_i(t),
\]

where $u_i(t)$ denotes active practice and maintenance.

The solution is

\[
m_i(t)
=
m_i(0)e^{-\lambda_i t}
+
\int_0^t
e^{-\lambda_i(t-\tau)}
u_i(\tau)
\, d\tau.
\]

When maintenance activity vanishes,

\[
u_i(t)=0,
\]

capability decays exponentially:

\[
m_i(t)
=
m_i(0)e^{-\lambda_i t}.
\]

Operational Fiber Death corresponds to

\[
m_i(t)
<
m_{\mathrm{crit}}
\]

for all capabilities required to reconstruct a historical
institutional configuration.

\section{C. Attractor Transition Criterion}

Consider the institutional potential

\[
V(x,t)
=
-
\langle P(t),F(x)\rangle.
\]

An attractor state $x^\ast$ satisfies

\[
\nabla V(x^\ast,t)=0.
\]

Local stability requires

\[
H(x^\ast,t)
=
\nabla^2 V(x^\ast,t)
\succ 0.
\]

Define the instability index

\[
I(t)
=
-\lambda_{\min}
\big(
H(x^\ast,t)
\big),
\]

where $\lambda_{\min}$ is the smallest eigenvalue.

Then

\[
I(t) > 0
\]

marks loss of local stability.

Institutional crises can therefore be interpreted as bifurcation
events in which

\[
\lambda_{\min}
\rightarrow 0.
\]

The transition period corresponds to a temporary expansion of
accessible trajectories before a new attractor basin forms.


\section{D. Opacity as Information Compression}

Let

\[
X
\]

denote the full institutional state and

\[
M=\pi(X)
\]

the publicly observable projection.

Opacity was defined as

\[
\Omega
=
H(X)-H(M).
\]

We may normalize this quantity by defining

\[
\omega
=
\frac{\Omega}{H(X)}
=
1-\frac{H(M)}{H(X)}.
\]

Thus

\[
0 \le \omega \le 1.
\]

A value near zero corresponds to institutional transparency.

A value near one corresponds to near-complete information loss under
projection.

The apparent agency threshold may therefore be written

\[
A(t)\,\omega(t)
>
\Theta,
\]

where $\Theta$ is the minimum alignment-opacity product required for
observers to attribute unified intention to the system.

\section{E. Functional Reconstruction Under Memory Loss}

Let

\[
S_{\mathrm{old}}
\]

denote a historical institution and

\[
S_{\mathrm{new}}
\]

a proposed replacement.

Define a functional map

\[
\Phi(S)
=
(f_1,f_2,\ldots,f_k),
\]

where each coordinate measures an institutional function.

Examples include

\[
f_1
=
\text{employment stability},
\qquad
f_2
=
\text{housing accessibility},
\qquad
f_3
=
\text{capital discipline}.
\]

Exact restoration requires

\[
S_{\mathrm{new}}
\cong
S_{\mathrm{old}}.
\]

Approximate repair requires only

\[
\|
\Phi(S_{\mathrm{new}})
-
\Phi(S_{\mathrm{old}})
\|
<
\epsilon.
\]

Theorem~3 follows immediately:

\[
S_{\mathrm{new}}
\not\cong
S_{\mathrm{old}}
\quad\text{but}\quad
\Phi(S_{\mathrm{new}})
\approx
\Phi(S_{\mathrm{old}}).
\]

The object of institutional repair is therefore preservation of function rather than preservation of form.

\newpage

% ---------------------------------------------------------------
\begin{thebibliography}{99}

% Formal and systems foundations
\bibitem{wiener1948}
Wiener, N. (1948).
\textit{Cybernetics: Or Control and Communication in the Animal and the Machine}.
MIT Press.

\bibitem{strogatz2015}
Strogatz, S. H. (2015).
\textit{Nonlinear Dynamics and Chaos}.
Westview Press.

\bibitem{simon1962}
Simon, H. A. (1962).
The Architecture of Complexity.
\textit{Proceedings of the American Philosophical Society}, 106(6), 467--482.

\bibitem{holland1992}
Holland, J. H. (1992).
\textit{Adaptation in Natural and Artificial Systems}.
MIT Press.

\bibitem{holland1995}
Holland, J. H. (1995).
\textit{Hidden Order: How Adaptation Builds Complexity}.
Addison-Wesley.

\bibitem{meadows2008}
Meadows, D. H. (2008).
\textit{Thinking in Systems: A Primer}.
Chelsea Green.

\bibitem{page2018}
Page, S. E. (2018).
\textit{The Model Thinker}.
Basic Books.

\bibitem{taleb2012}
Taleb, N. N. (2012).
\textit{Antifragile: Things That Gain from Disorder}.
Random House.

% Institutional theory and organizational memory
\bibitem{douglassnorth1990}
North, D. C. (1990).
\textit{Institutions, Institutional Change and Economic Performance}.
Cambridge University Press.

\bibitem{nelsonwinter1982}
Nelson, R. R., \& Winter, S. G. (1982).
\textit{An Evolutionary Theory of Economic Change}.
Harvard University Press.

\bibitem{cyertmarch1963}
Cyert, R. M., \& March, J. G. (1963).
\textit{A Behavioral Theory of the Firm}.
Prentice Hall.

\bibitem{marcholsen1989}
March, J. G., \& Olsen, J. P. (1989).
\textit{Rediscovering Institutions: The Organizational Basis of Politics}.
Free Press.

\bibitem{march1991}
March, J. G. (1991).
Exploration and Exploitation in Organizational Learning.
\textit{Organization Science}, 2(1), 71--87.

\bibitem{weick1995}
Weick, K. E. (1995).
\textit{Sensemaking in Organizations}.
Sage.

\bibitem{teece1997}
Teece, D. J., Pisano, G., \& Shuen, A. (1997).
Dynamic Capabilities and Strategic Management.
\textit{Strategic Management Journal}, 18(7), 509--533.

\bibitem{williamson1985}
Williamson, O. E. (1985).
\textit{The Economic Institutions of Capitalism}.
Free Press.

\bibitem{ostrom1990}
Ostrom, E. (1990).
\textit{Governing the Commons: The Evolution of Institutions for Collective Action}.
Cambridge University Press.

% Path dependence and lock-in
\bibitem{arthur1989}
Arthur, W. B. (1989).
Competing Technologies, Increasing Returns, and Lock-In by Historical Events.
\textit{Economic Journal}, 99(394), 116--131.

\bibitem{arthur1994}
Arthur, W. B. (1994).
\textit{Increasing Returns and Path Dependence in the Economy}.
University of Michigan Press.

\bibitem{arrow1962}
Arrow, K. J. (1962).
The Economic Implications of Learning by Doing.
\textit{Review of Economic Studies}, 29(3), 155--173.

% Paradigm change and scientific revolutions
\bibitem{kuhn1962}
Kuhn, T. S. (1962).
\textit{The Structure of Scientific Revolutions}.
University of Chicago Press.

% Social embeddedness and power
\bibitem{polanyi1944}
Polanyi, K. (1944).
\textit{The Great Transformation: The Political and Economic Origins of Our Time}.
Beacon Press.

\bibitem{granovetter1985}
Granovetter, M. (1985).
Economic Action and Social Structure: The Problem of Embeddedness.
\textit{American Journal of Sociology}, 91(3), 481--510.

\bibitem{sen1999}
Sen, A. (1999).
\textit{Development as Freedom}.
Oxford University Press.

% Knowledge and distributed cognition
\bibitem{hayek1945}
Hayek, F. A. (1945).
The Use of Knowledge in Society.
\textit{American Economic Review}, 35(4), 519--530.

% Political economy of the neoliberal transition
\bibitem{powell1971}
Powell, L. F. (1971).
\textit{Attack on the American Free Enterprise System}.
Confidential Memorandum to the U.S.\ Chamber of Commerce.

\bibitem{harvey2005}
Harvey, D. (2005).
\textit{A Brief History of Neoliberalism}.
Oxford University Press.

\bibitem{mirowski2013}
Mirowski, P. (2013).
\textit{Never Let a Serious Crisis Go to Waste: How Neoliberalism Survived the Financial Meltdown}.
Verso.

\bibitem{streeck2016}
Streeck, W. (2016).
\textit{How Will Capitalism End? Essays on a Failing System}.
Verso.

\bibitem{krippner2011}
Krippner, G. R. (2011).
\textit{Capitalizing on Crisis: The Political Origins of the Rise of Finance}.
Harvard University Press.

\bibitem{dumenil2004}
Dum\'enil, G., \& L\'evy, D. (2004).
\textit{Capital Resurgent: Roots of the Neoliberal Revolution}.
Harvard University Press.

\bibitem{piketty2014}
Piketty, T. (2014).
\textit{Capital in the Twenty-First Century}.
Harvard University Press.

\bibitem{minsky1986}
Minsky, H. P. (1986).
\textit{Stabilizing an Unstable Economy}.
Yale University Press.

\bibitem{galbraith1967}
Galbraith, J. K. (1967).
\textit{The New Industrial State}.
Houghton Mifflin.

\bibitem{rodrik2011}
Rodrik, D. (2011).
\textit{The Globalization Paradox: Democracy and the Future of the World Economy}.
W.\ W.\ Norton.

\bibitem{stiglitz2002}
Stiglitz, J. E. (2002).
\textit{Globalization and Its Discontents}.
W.\ W.\ Norton.

% Cliodynamics and long-run cycles
\bibitem{turchin2016}
Turchin, P. (2016).
\textit{Ages of Discord: A Structural-Demographic Analysis of American History}.
Beresta Books.

\end{thebibliography}

% ---------------------------------------------------------------
\end{document}
